\ifdefined\gkver\else\def\gkver{student}\fi
\documentclass[\gkver]{gaokaozhenti}
\begin{document}

\chapter{2007年江苏}

\begin{enumerate}
\item
%% number: 1
%% paperName: 2007年江苏高考第 1 题 3 分
%% typeId: 1
%% type: 单选题
%% chapter: 气体性质
%% point: 分子力
%% method: 无
%% score: 3
%% degree: 650
%% duplicateId: 0
%% body:
分子动理论较好地解释了物质的宏观热力学性质。据此可判断下列说法中错误的是\xzanswer{B}

\fourchoices[answer=B]
{显微镜下观察到墨水中的小炭粒在不停的作无规则运动，这反映了液体分子运动的无规则性}
{分子间的相互作用力随着分子间距离的增大，一定先减小后增大}
{分子势能随着分子间距离的增大，可能先减小后增大}
{在真空、高温条件下，可以利用分子扩散向半导体材料掺入其它元素}

%% answer: B
%% memo:
\memoanswer{
墨水中的小炭粒做布朗运动，反映了液体分子做无规则运动，故 A 正确；分子间的相互作用力表现为斥力时，分子间距离越大，分子力越小；当分子间的作用力表现为引力，分子间的距离由最小开始变大时，分子间作用力先增大，后减小，故 B 错误；分子势能随着分子间距离的增大，可能先减小后增大，故 C 正确；由于分子在做永不停息的无规则运动，高温情况下分子能很快地扩散到其他材料中，故 D 正确。
}

\item
%% number: 2
%% paperName: 2007年江苏高考第 2 题 3 分
%% typeId: 1
%% type: 单选题
%% chapter: 原子和原子核
%% point: 人工核反应
%% method: 无
%% score: 3
%% degree: 650
%% duplicateId: 0
%% body:
2006 年美国和俄罗斯的科学家利用回旋加速器，通过 \ce{^{48}_{20}Ca} 轰击 \ce{^{249}_{98}Cf} 发生核反应，成功合成了第 118 号元素，这是迄今为止门捷列夫元素周期表中原子序数最大的元素，实验表明，该元素的原子核先放出 3 个相同的粒子 $x$，再连续经过 3 次 $\alpha$ 衰变后，变成质量为 282 的第 112 号元素的原子核，则上述过程中的粒子 $x$ 是\xzanswer{A}

\fourchoices[answer=A]
{中子}
{质子}
{电子}
{$\alpha$ 粒子}

%% answer: A
%% memo:
\memoanswer{
由 \ce{^{48}_{20}Ca} 与 \ce{^{249}_{98}Cf} 合成的第 118 号元素质量数是 297；发生 3 次 $\alpha$ 衰变，质量数少 12、核电荷数少了 6，可见先放出的三个粒子电荷数是 0，质量数是 1，是中子，故 A 正确。
}

\item
%% number: 3
%% paperName: 2007年江苏高考第 3 题 3 分
%% typeId: 1
%% type: 单选题
%% chapter: 光学
%% point: 光的电磁说
%% method: 无
%% score: 3
%% degree: 650
%% duplicateId: 0
%% body:
光的偏振现象说明光是横波，下列现象中不能反映光的偏振特性的是\xzanswer{D}

\fourchoices[answer=D]
{一束自然光相继通过两个偏振片，以光束为轴旋转其中一个偏振片，透射光的强度发生变化}
{一束自然光入射到两种介质的分界面上，当反射光与折射光线之间的夹角恰好是 $90^{\circ}$ 时，反射光是偏振光}
{日落时分，拍摄水面下的景物，在照相机镜头前装上偏振光片可以使景象更清晰}
{通过手指间的缝隙观察日光灯，可以看到彩色条纹}

%% answer: D
%% memo:
\memoanswer{
通过手指的缝隙观察日光灯，看到彩色条纹，是光的衍射现象，故 D 错误。
}

\item
%% number: 4
%% paperName: 2007年江苏高考第 4 题 3 分
%% typeId: 1
%% type: 单选题
%% chapter: 原子和原子核
%% point: 玻尔模型
%% method: 无
%% score: 3
%% degree: 650
%% duplicateId: 0
%% body:
$\mu$ 子与氢原子核（质子）构成的原子称为 $\mu$ 氢原子（hydrogen muon atom），它在原子核物理的研究中有重要作用。图为 $\mu$ 氢原子的能级示意图。假定光子能量为 $E$ 的一束光照射容器中大量处于 $n=2$ 能级的 $\mu$ 氢原子，$\mu$ 氢原子吸收光子后，发出频率为 $\nu_1$、$\nu_2$、$\nu_3$、$\nu_4$、$\nu_5$ 和 $\nu_6$ 的光，且频率依次增大，则 $E$ 等于\xzanswer{C}

\onepicture[w=8.0cm]{\includesvg{figs/04}}

\fourchoices[answer=C]
{$h(\nu_3-\nu_1)$}
{$h(\nu_5+\nu_6)$}
{$h\nu_3$}
{$h\nu_4$}

%% answer: C
%% memo:
\memoanswer{
$\mu$ 氢原子吸收某频率为 $\nu$ 的光子后，跃迁到第 $n$ 能级，由于能发出 6 种光子，说明 $n=4$，从 $n=4$ 跃迁到各个较低能级上放出的光子的频率如图所示，可知光子的能量是 $E=h\nu_3$，故 C 正确。

\onepicture[w=8.5cm]{figs/04_an.jpg}
}

\item
%% number: 5
%% paperName: 2007年江苏高考第 5 题 3 分
%% typeId: 1
%% type: 单选题
%% chapter: 机械波
%% point: 波的多解性
%% method: 无
%% score: 3
%% degree: 650
%% duplicateId: 0
%% body:
如图所示，实线和虚线分别为某种波在 $t$ 时刻和 $t+\Delta t$ 时刻的波形曲线。B 和 C 是横坐标分别为 $d$ 和 $3d$ 的两个质点，下列说法中正确的是\xzanswer{C}

\onepicture[w=11.0cm]{\includesvg{figs/05}}

\fourchoices[answer=C]
{任一时刻，如果质点 B 向上运动，则质点 C 一定向下运动}
{任一时刻，如果质点 B 速度为零，则质点 C 的速度也为零}
{如果波是向右传播的，则波的周期可能为 $\frac{6}{7}\Delta t$}
{如果波是向左传播的，则波的周期可能为 $\frac{6}{13}\Delta t$}

%% answer: C
%% memo:
\memoanswer{
把实线波形向左平移 $d$，并利用波的传播方向与质点振动方向之间的关系，能判断出 B、C 的振动方向可能都向上，故 A 错误；把波向左平移 $\frac{d}{4}$，质点 B 的速度为零，质点 C 的速度不为零，故 B 错误；波向右传播时，向前传播的距离是 $x=\frac{d}{2}+3dn$（$n=0,1,2,\ldots$），由 $v=\frac{\lambda}{T}=\frac{\Delta x}{\Delta t}$，$T=\frac{\lambda\Delta t}{\Delta x}=\frac{3d\Delta t}{\frac{d}{2}+3dn}$ 可知，波的周期可能为 $\frac{6}{7}\Delta t$，故 C 正确；如果波向左传播，向前传播的距离是 $x=\frac{5d}{2}+3dn$（$n=0,1,2,\ldots$），由 $v=\frac{\lambda}{T}=\frac{\Delta x}{\Delta t}$，$T=\frac{\lambda\Delta t}{\Delta x}=\frac{3d\Delta t}{\frac{5d}{2}+3dn}$ 可知，波的周期不可能是 $\frac{6}{13}\Delta t$，故 D 错误。
}

\item
%% number: 6
%% paperName: 2007年江苏高考第 6 题 3 分
%% typeId: 1
%% type: 单选题
%% chapter: 运动和力
%% point: 连接体
%% method: 临界
%% score: 3
%% degree: 450
%% duplicateId: 0
%% body:
如图所示，光滑水平面上放置质量分别为 $m$ 和 $2m$ 的四个木块，其中两个质量为 $m$ 的木块间用一不可伸长的轻绳相连，木块间的最大静摩擦力是 $\mu mg$。现用水平拉力 $F$ 拉其中一个质量为 $2m$ 的木块，使四个木块以同一加速度运动，则轻绳对 $m$ 的最大拉力为\xzanswer{B}

\onepicture[w=10.0cm]{\includesvg{figs/06}}

\fourchoices[answer=B]
{$\frac{3}{5}\mu mg$}
{$\frac{3}{4}\mu mg$}
{$\frac{3}{2}\mu mg$}
{$3\mu mg$}

%% answer: B
%% memo:
\memoanswer{
对四个木块组成的整体受力分析，设整体加速度为 $a$，水平方向受拉力 $F$ 的作用，$F=6ma$；如果对左边的两个木块组成的整体受力分析，水平方向只受拉力 $T$，$T=3ma=\frac{F}{2}$；对左边质量 $2m$ 的木块受力分析，水平方向只受静摩擦力作用，$f=2ma$；对除右边质量 $2m$ 以外的三个木块组成的整体受力分析，水平方向只受静摩擦力作用，由 $f=4ma$，所以这时整体受到的静摩擦力最大，$f=\mu mg$，$a=\frac{\mu g}{4}$，$T=\frac{3}{4}\mu mg$，故 B 正确。
}

\item
%% number: 7
%% paperName: 2007年江苏高考第 7 题 4 分
%% typeId: 2
%% type: 多选题
%% chapter: 光学
%% point: 波粒二象性
%% method: 无
%% score: 4
%% degree: 650
%% duplicateId: 0
%% body:
现代物理学认为，光和实物粒子都具有波粒二象性。下列事实中突出体现波动性的是\xzanswer{BD}

\fourchoices[answer=BD]
{一定频率的光照射到锌板上，光的强度越大，单位时间内锌板上发射的光电子就越多}
{肥皂液是无色的，吹出的肥皂泡却是彩色的}
{质量为 $10^{-3}\Ukg$、速度为 $10^{-2}\Ums$ 的小球，其德布罗意波长约为 $10^{-23}\Um$，不过我们能清晰地观测到小球运动的轨迹}
{人们常利用热中子研究晶体的结构，因为热中子的德布罗意波长与晶体中原子间距大致相同}

%% answer: BD
%% memo:
\memoanswer{
光子照到锌板上，发生光电效应，说明光有粒子性，故 A 错误；白光在肥皂泡上发生薄膜干涉时，会出现彩色条纹，光的干涉现象说明了光有波动性，故 B 正确；由于实物的波长很小，波动性不明显，表现为粒子性，所以能看到小球的轨迹，故 C 错误；用热中子研究晶体结构，其实是通过中子的衍射来“观察”晶体的，是利用中子的波动性，故 D 正确。
}

\item
%% number: 8
%% paperName: 2007年江苏高考第 8 题 4 分
%% typeId: 2
%% type: 多选题
%% chapter: 光学
%% point: 光的电磁说
%% method: 无
%% score: 4
%% degree: 450
%% duplicateId: 0
%% body:
2006 年度诺贝尔物理学奖授予了两名美国科学家，以表彰他们发现了宇宙微波背景辐射的黑体谱形状及其温度在不同方向上的微小变化。他们的出色工作被誉为是宇宙学研究进入精密科学时代的起点，下列与宇宙微波背景辐射黑体谱相关的说法中正确的是\xzanswer{ACD}

\fourchoices[answer=ACD]
{微波是指波长在 $10^{-3}\Um$ 到 $10\Um$ 之间的电磁波}
{微波和声波一样都只能在介质中传播}
{黑体的热辐射实际上是电磁辐射}
{普朗克在研究黑体的热辐射问题中提出了能量子假说}

%% answer: ACD
%% memo:
\memoanswer{
微波的波长范围是 $1\Umm$ 到 $10\Um$，故 A 正确；微波能在真空中传播，故 B 错误；黑体热辐射实际上是电磁辐射，故 C 正确；普朗克在研究黑体的热辐射问题时提出了能量子假说，故 D 正确。
}

\item
%% number: 9
%% paperName: 2007年江苏高考第 9 题 4 分
%% typeId: 2
%% type: 多选题
%% chapter: 电路
%% point: 电路实验
%% method: 无
%% score: 4
%% degree: 450
%% duplicateId: 0
%% body:
某同学欲采用如图所示的电路完成相关实验。图中电流表的量程为 $0.6\UA$，内阻约 $0.1\UO$；电压表的量程为 $3\UV$，内阻约 $6\UkO$；G 为小量程电流表；电源电动势约 $3\UV$，内阻较小，下列电路中正确的是\xzanswer{AB}

\fourchoices[answer=AB, ispicture=true, h=2.6cm]
{figs/09a.png}
{figs/09b.png}
{figs/09c.png}
{figs/09d.png}

%% answer: AB
%% memo:
\memoanswer{
电路 A、C 测量的是相对较小的电阻，用安培表外接法是恰当的，A 图中滑动变阻器用限流接法也符合实验要求，故 A 正确；C 图中电路要描绘小灯泡的伏安曲线，需要在从零开始的较大范围内调节电压，故 C 错误；B 图是测定电源电动势和内电阻的电路，是正确的，故 B 正确；D 图中变阻箱与滑动变阻器的位置颠倒，故 D 错误。
}

\item
%% number: 10
%% paperName: 2007年江苏高考第 10 题 4 分
%% typeId: 2
%% type: 多选题
%% chapter: 曲线运动
%% point: 人造地球卫星
%% method: 无
%% score: 4
%% degree: 450
%% duplicateId: 0
%% body:
假设太阳系中天体的密度不变，天体直径和天体之间距离都缩小到原来的一半，地球绕太阳公转近似为匀速圆周运动，则下列物理量变化正确的是\xzanswer{BC}

\fourchoices[answer=BC]
{地球的向心力变为缩小前的一半}
{地球的向心力变为缩小前的 $\frac{1}{16}$}
{地球绕太阳公转周期与缩小前的相同}
{地球绕太阳公转周期变为缩小前的一半}

%% answer: BC
%% memo:
\memoanswer{
地球质量 $m=\rho_1\cdot\frac{4}{3}\pi R_1^3$，太阳质量 $M=\rho_2\cdot\frac{4}{3}\pi R_2^3$，由万有引力定律得 $F_{\text{万}}=G\frac{Mm}{r^2}=\frac{16\pi^2G\rho_1\rho_2R_1^3R_2^3}{9r^2}$，地球受到的万有引力等于向心力，当天体直径和距离 $r$ 缩小一半时，向心力变为原来的 $\frac{1}{16}$，故 B 正确 A 错误；又 $G\frac{Mm}{r^2}=m\frac{4\pi^2}{T^2}r$，解得 $T^2=\frac{4\pi^2r^3}{GM}=\frac{4\pi^2r^3}{G\rho_2\cdot\frac{4}{3}\pi R_2^3}$，当天体直径和距离 $r$ 缩小一半时，$T$ 不变，故 C 正确 D 错误。
}

\item
%% number: 11
%% paperName: 2007年江苏高考第 11 题 4 分
%% typeId: 2
%% type: 多选题
%% chapter: 气体性质
%% point: 热力学第一定律
%% method: 无
%% score: 4
%% degree: 650
%% duplicateId: 0
%% body:
如图所示，绝热气缸中间用固定栓将可无摩擦移动的导热隔板固定，隔板质量不计，左右两室分别充有一定量的氢气和氧气（视为理想气体）。初始时，两室气体的温度相等，氢气的压强大于氧气的压强，松开固定栓直至系统重新达到平衡，下列说法中正确的是\xzanswer{CD}

\onepicture[w=8.0cm]{\includesvg{figs/11}}

\fourchoices[answer=CD]
{初始时氢分子的平均动能大于氧分子的平均动能}
{系统重新达到平衡时，氢气的内能比初始时的小}
{松开固定栓直至系统重新达到平衡的过程中有热量从氧气传递到氢气}
{松开固定栓直至系统重新达到平衡的过程中，氧气的内能先增大后减小}

%% answer: CD
%% memo:
\memoanswer{
温度是分子平均动能大小的标志，初始时，氢气和氧气的温度相同，分子平均动能相同，故 A 错误；气体与气缸外部没有热传递和其他能量交换，气缸内气体内能不会变化，故 B 错误；在刚松开固定栓时，氢气膨胀，对氧气做功，氢气内能减小，氢气温度降低，氧气内能增加，氧气温度升高；由于隔板是导热的，热量从氧气传递到氢气，直至温度相同，氧气的内能又减小，故 CD 正确。
}

\item
%% number: 12
%% paperName: 2007年江苏高考第 12 题 9 分
%% typeId: 4
%% type: 实验题
%% chapter: 电路
%% point: 电路实验
%% method: 无
%% score: 9
%% degree: 450
%% duplicateId: 0
%% body:
要描绘某电学元件（最大电流不超过 $6\UmA$，最大电压不超过 $7\UV$）的伏安特性曲线，设计电路如图，图中定值电阻 $R$ 为 $1\UkO$，用于限流；电流表量程为 $10\UmA$，内阻约为 $5\UO$；电压表（未画出）量程为 $10\UV$，内阻约为 $10\UkO$；电源电动势 $E$ 为 $12\UV$，内阻不计。

\onepicture[w=6.5cm]{\includesvg{figs/12a}}

\begin{enumerate}
	\item
	实验时有两个滑动变阻器可供选择：
	
	a、阻值 0 到 $200\UO$，额定电流 $0.3\UA$；
	
	b、阻值 0 到 $20\UO$，额定电流 $0.5\UA$。
	
	本实验应选的滑动变阻器是\tkanswer[1.5cm]{}（填“a”或“b”）。
	
	\item
	正确接线后，测得数据如下表
	
	\begin{table}[htbp]
	\centering
	\begin{tblr}{|c|c|c|c|c|c|c|c|c|c|c|}
	\hline
	 & 1 & 2 & 3 & 4 & 5 & 6 & 7 & 8 & 9 & 10 \\
	\hline
	$U$（$\UV$） & 0.00 & 3.00 & 6.00 & 6.16 & 6.28 & 6.32 & 6.36 & 6.38 & 6.39 & 6.40 \\
	\hline
	$I$（$\UmA$） & 0.00 & 0.00 & 0.00 & 0.06 & 0.50 & 1.00 & 2.00 & 3.00 & 4.00 & 5.50 \\
	\hline
	\end{tblr}
	\end{table}
	
	①根据以上数据，电压表是并联在 M 与\tkanswer[1.2cm]{}之间的（填“O”或“P”）；
	
	②根据以上数据，画出该元件的伏安特性曲线。
	
	\onepicture[w=6.5cm]{\includesvg{figs/12b}}
	
	\item
	画出待测元件两端电压 $U_{\mathrm{MO}}$ 随 MN 间电压 $U_{\mathrm{MN}}$ 变化的示意图。（无需数值）
	
	\onepicture[w=6.5cm]{\includesvg{figs/12c}}
\end{enumerate}

%% answer: 
\jdanswer{
\begin{enumerate}
	\item
	a
	
	\item
	①P；②如图所示
	
	\item
	如图所示
\end{enumerate}
}
%% memo:
\memoanswer{
（1）选择滑动变阻器首先要考虑安全，如果选用 b 滑动变阻器，滑动变阻器中的电流将超过 $0.5\UA$ 的额定电流，因此不合适，故选 a。

（2）从表格数据可以看出，元件电阻在 $1000\UO$ 以上，相对电表内阻是大电阻，电压表分流作用明显，应接 MP 之间；可以计算出 MN 之间的电压分别是 $0.00\UV$、$3.00\UV$、$6.00\UV$、$6.22\UV$、$6.78\UV$、$7.32\UV$、$8.36\UV$、$9.38\UV$、$10.39\UV$、$11.90\UV$，通过粗略描点可以得到如图所示的图象。

\onepicture[w=7.0cm]{\includesvg{figs/12_an1}}

（3）画出待测元件两端电压 $U_{\mathrm{MO}}$ 随 $U_{\mathrm{MN}}$ 变化的示意图如图所示。

\onepicture[w=7.0cm]{\includesvg{figs/12_an2}}
}

\item
%% number: 13
%% paperName: 2007年江苏高考第 13 题 13 分
%% typeId: 4
%% type: 实验题
%% chapter: 运动和力
%% point: 验证牛顿第二定律
%% method: 无
%% score: 13
%% degree: 450
%% duplicateId: 0
%% body:
如（\ref{2007江苏13a}）图，质量为 $M$ 的滑块 A 放在气垫导轨 B 上，C 为位移传感器，它能将滑块 A 到传感器 C 的距离数据实时传送到计算机上，经计算机处理后在屏幕上显示滑块 A 的位移-时间（$s-t$）图象和速率-时间（$v-t$）图象。整个装置置于高度可调节的斜面上，斜面的长度为 $l$、高度为 $h$。（取重力加速度 $g=9.8\Umsq$，结果可保留一位有效数字）

\begin{enumerate}
	\item
	现给滑块 A 一沿气垫导轨向上的初速度，A 的 $v-t$ 图线如（\ref{2007江苏13b}）图所示。从图线可得滑块 A 下滑时的加速度 $a=$\tkanswer[1.5cm]{}$\Umsq$，摩擦力对滑块 A 运动的影响\tkanswer[3.0cm]{}。（填“明显，不可忽略”或“不明显，可忽略”）
	
	\item
	此装置还可用来验证牛顿第二定律。实验时通过改变\tkanswer[2.5cm]{}，可验证质量一定时，加速度与力成正比的关系；实验时通过改变\tkanswer[3.0cm]{}，可验证力一定时，加速度与质量成反比的关系。
	
	\item
	将气垫导轨换成滑板，滑块 A 换成滑块 $A'$，给滑块 $A'$ 一沿滑板向上的初速度，$A'$ 的 $s-t$ 图线如（\ref{2007江苏13c}）图。图线不对称是由于\tkanswer[3.0cm]{}造成的，通过图线可求得滑板的倾角 $\theta=$\tkanswer[2.5cm]{}（用反三角函数表示），滑块与滑板间的动摩擦因数 $\mu=$\tkanswer[1.5cm]{}。
\end{enumerate}

\threepicture[h=3.4cm, align=right, labelA=2007江苏13a, labelB=2007江苏13b, labelC=2007江苏13c]
{figs/13a.png}
{figs/13b.png}
{figs/13c.png}

%% answer: 
\jdanswer{
\begin{enumerate}
	\item
	6；不明显，可忽略
	
	\item
	斜面高度 $h$；滑块 A 的质量 $M$ 及斜面的高度 $h$，且使 $Mh$ 不变
	
	\item
	滑动摩擦力；$\arcsin 0.6$（$\arcsin 0.57\sim\arcsin 0.64$ 都算对）；0.3（0.2$\sim$0.4 都算对）
\end{enumerate}
}
%% memo:
\memoanswer{
（1）从图象可以看出，滑块上滑和下滑过程中的加速度基本相等，所以摩擦力对滑块的运动影响不明显，可以忽略。根据加速度的定义式可以得出 $a=\frac{\Delta v}{\Delta t}=\frac{2.9}{0.48}\Umsq=6\Umsq$。

（2）牛顿第二定律研究的是加速度与合外力和质量的关系。当质量一定时，可以改变力的大小，当斜面高度不同时，滑块受到的力不同，可以探究加速度与合外力的关系。由于滑块下滑加速的力是由重力沿斜面向下的分力提供，所以要保证向下的分力不变，应该使 $Mg\cdot\frac{h}{l}$ 不变，所以应该调节滑块 A 的质量 $M$ 及斜面的高度 $h$，且使 $Mh$ 不变。

（3）滑板与滑块间的滑动摩擦力比较大，导致图象成抛物线形。从图上可以读出，物体沿斜面向上滑动时，加速度 $a_1=g\sin\theta+\mu g\cos\theta$，由运动学公式得 $s_1=\frac{1}{2}a_1t_1^2$，$s_1=0.81\Um-0.20\Um$，$t_1=0.4\Us$，物体沿斜面向下滑动时，加速度 $a_2=g\sin\theta-\mu g\cos\theta$，由运动学公式得 $s_2=\frac{1}{2}a_2t_2^2$，$s_2=0.84\Um-0.20\Um$，$t_2=0.6\Us$，由上面的关系可以解出 $\theta=\arcsin 0.6$，$\mu=0.3$。
}

\item
%% number: 14
%% paperName: 2007年江苏高考第 14 题 14 分
%% typeId: 6
%% type: 计算题
%% chapter: 光学
%% point: 光的折射
%% method: 无
%% score: 14
%% degree: 450
%% duplicateId: 0
%% body:
如图所示，巡查员站立于一空的贮液池边，检查池角处出液口的安全情况。已知池宽为 $L$，照明灯到池底的距离为 $H$。若保持照明光束方向不变，向贮液池中注入某种液体，当液面高为 $H/2$ 时，池底的光斑距离出液口 $L/4$。

\begin{enumerate}
	\item
	试求当液面高为 $\frac{2}{3}H$ 时，池底的光斑到出液口的距离 $x$。
	
	\item
	控制出液口缓慢地排出液体，使液面以 $v_0$ 的速率匀速下降，试求池底的光斑移动的速率 $v_x$。
\end{enumerate}

\onepicture[w=8.0cm, align=right]{\includesvg{figs/14}}

%% answer: 
\jdanswer{
\begin{enumerate}
	\item
	$x=\frac{L}{3}$
	
	\item
	$v_x=\frac{L}{2H}v_0$
\end{enumerate}
}
%% memo:
\memoanswer{
（1）几何关系如图所示，由几何关系知 $\frac{x+l}{h}=\frac{L}{H}$ ①，$h=\frac{H}{2}$ 时，$x=\frac{L}{4}$，故 $\tan\gamma=\frac{l}{h}=\frac{\frac{L}{4}}{\frac{H}{2}}=\frac{L}{2H}$ ②，因为液面高度变化，折射角 $\gamma$ 不变，故由①式得 $\frac{x+l}{h}=\frac{x}{h}+\frac{l}{h}=\frac{L}{H}$ ③，②代入③得 $x=\frac{L}{2H}h$ ④，将 $h=\frac{2}{3}H$ 代入④中得 $x=\frac{L}{3}$。

（2）因为 $\Delta x=\frac{L}{2H}\Delta h$，所以 $v_x=\frac{L}{2H}\cdot v_0$。

\onepicture[w=8.0cm]{figs/14_an.jpg}
}

\item
%% number: 15
%% paperName: 2007年江苏高考第 15 题 14 分
%% typeId: 6
%% type: 计算题
%% chapter: 运动和力
%% point: 牛顿第二定律
%% method: 无
%% score: 14
%% degree: 650
%% duplicateId: 0
%% body:
直升机沿水平方向匀速飞往水源取水灭火，悬挂着 $m=500\Ukg$ 空箱的悬索与竖直方向的夹角 $\theta_1=45^{\circ}$。直升机取水后飞往火场，加速度沿水平方向，大小稳定在 $a=1.5\Umsq$ 时，悬索与竖直方向的夹角 $\theta_2=14^{\circ}$。如果空气阻力大小不变，且忽略悬索的质量，试求水箱中水的质量 $M$。（取重力加速度 $g=10\Umsq$；$\sin 14^{\circ}=0.242$；$\cos 14^{\circ}=0.970$）

\onepicture[w=10.0cm, align=right]{\includesvg{figs/15}}

%% answer: 
\jdanswer{
$M=4.5\times10^3\Ukg$
}
%% memo:
\memoanswer{
直升机取水，水箱受力平衡：
$T_1\sin\theta_1-f=0$ ①，
$T_1\cos\theta_1-mg=0$ ②，
由①②解得 $f=mg\tan\theta_1$ ③。

直升机返回，由牛顿第二定律得：
$T_2\sin\theta_2-f=(M+m)a$ ④，
$T_2\cos\theta_2-(M+m)g=0$ ⑤，
联立解得水箱中水的质量为 $M=4.5\times10^3\Ukg$。
}

\item
%% number: 16
%% paperName: 2007年江苏高考第 16 题 15 分
%% typeId: 6
%% type: 计算题
%% chapter: 电场
%% point: 带电粒子的平衡
%% method: 无
%% score: 15
%% degree: 450
%% duplicateId: 0
%% body:
如图所示，带电量分别为 $4q$ 和 $-q$ 的小球 A、B 固定在水平放置的光滑绝缘细杆上，相距为 $d$。若杆上套一带电小环 C，带电体 A、B 和 C 均可视为点电荷。

\begin{enumerate}
	\item
	求小环 C 的平衡位置。
	
	\item
	若小环 C 带电量为 $q$，将小环拉离平衡位置一小位移 $x$（$|x|\ll d$）后静止释放，试判断小环 C 能否回到平衡位置。（回答“能”或“不能”即可）
	
	\item
	若小环 C 带电量为 $-q$，将小环拉离平衡位置一小位移 $x$（$|x|\ll d$）后静止释放，试证明小环 C 将作简谐运动。（提示：当 $\alpha\ll1$ 时，则 $\frac{1}{(1+\alpha)^n}\approx1-n\alpha$）
\end{enumerate}

\onepicture[w=10.0cm, align=right]{\includesvg{figs/16}}

%% answer: 
\jdanswer{
\begin{enumerate}
	\item
	平衡位置为 $l=d$
	
	\item
	不能
	
	\item
	利用近似关系化简得 $F_C=-\frac{kq^2}{d^3}x$，所以小环 C 将做简谐运动
\end{enumerate}
}
%% memo:
\memoanswer{
（1）设 C 在 AB 连线的延长线上距离 B 为 $l$ 处达到平衡，带电量为 $Q$。由库仑定律得 $F=k\frac{qQ}{r^2}$，由平衡条件得 $F_C=\frac{4kqQ}{(d+l)^2}+\frac{-kqQ}{l^2}=0$，解得 $l_1=-\frac{1}{3}d$（舍去），$l_2=d$。所以平衡位置为 $l=d$。

（2）不能。

（3）环 C 带电 $-q$，平衡位置不变，拉离平衡位置一小位移 $x$ 后，C 受力为 $F_C=\frac{-4kq^2}{(2d+x)^2}+\frac{kq^2}{(d+x)^2}$，利用近似关系化简得 $F_C=-\frac{kq^2}{d^3}x$，所以小环 C 将做简谐运动。
}

\item
%% number: 17
%% paperName: 2007年江苏高考第 17 题 16 分
%% typeId: 6
%% type: 计算题
%% chapter: 磁场
%% point: 带电粒子在磁场中的运动
%% method: 无
%% score: 16
%% degree: 450
%% duplicateId: 0
%% body:
磁谱仪是测量 $\alpha$ 能谱的重要仪器。磁谱仪的工作原理如图所示，放射源 S 发出质量为 $m$、电量为 $q$ 的 $\alpha$ 粒子沿垂直磁场方向进入磁感应强度为 $B$ 的匀强磁场，被限束光栏 Q 限制在 $2\varphi$ 的小角度内，$\alpha$ 粒子经磁场偏转后打到与束光栏平行的感光片 P 上。（重力影响不计）

\begin{enumerate}
	\item
	若能量在 $E\sim E+\Delta E$（$\Delta E>0$，且 $\Delta E\ll E$）范围内的 $\alpha$ 粒子均垂直于限束光栏的方向进入磁场。试求这些 $\alpha$ 粒子打在胶片上的范围 $\Delta x_1$。
	
	\item
	实际上，限束光栏有一定的宽度，$\alpha$ 粒子将在 $2\varphi$ 角内进入磁场。试求能量均为 $E$ 的 $\alpha$ 粒子打到感光胶片上的范围 $\Delta x_2$。
\end{enumerate}

\onepicture[w=8.0cm, align=right]{\includesvg{figs/17}}

%% answer: 
\jdanswer{
\begin{enumerate}
	\item
	$\Delta x_1\approx\frac{\sqrt{2mE}}{qBE}\Delta E$
	
	\item
	$\Delta x_2=\frac{4\sqrt{2mE}}{qB}\sin^2\frac{\varphi}{2}$
\end{enumerate}
}
%% memo:
\memoanswer{
（1）设 $\alpha$ 粒子以速度 $v$ 进入磁场，打在胶片上的位置距 S 的距离为 $x$，由洛伦兹力提供向心力得 $qvB=m\frac{v^2}{R}$ ①，$\alpha$ 粒子的动能为 $E=\frac{1}{2}mv^2$ ②，又 $x=2R$ ③，由①②③式可得 $x=\frac{2\sqrt{2mE}}{qB}$ ④。

\onepicture[w=7.0cm]{figs/17_an1.jpg}

由④可得 $\Delta x_1=\frac{2\sqrt{2m(E+\Delta E)}}{qB}-\frac{2\sqrt{2mE}}{qB}$，化简可得 $\Delta x_1\approx\frac{\sqrt{2mE}}{qBE}\Delta E$。

（2）动能为 $E$ 的 $\alpha$ 粒子沿 $\pm\varphi$ 角入射，轨迹半径相同，设为 $R$，由洛伦兹力提供向心力得 $qvB=m\frac{v^2}{R}$，$\alpha$ 粒子的动能为 $E=\frac{1}{2}mv^2$。

\onepicture[w=7.0cm]{figs/17_an2.jpg}

由几何关系得 $\Delta x_2=2R-2R\cos\varphi=\frac{4\sqrt{2mE}}{qB}\sin^2\frac{\varphi}{2}$。
}

\item
%% number: 18
%% paperName: 2007年江苏高考第 18 题 16 分
%% typeId: 6
%% type: 计算题
%% chapter: 电磁感应
%% point: 电磁感应能量分析
%% method: 无
%% score: 16
%% degree: 450
%% duplicateId: 0
%% body:
如图所示，空间等间距分布着水平方向的条形匀强磁场，竖直方向磁场区域足够长，磁感应强度 $B=1\UT$，每一条形磁场区域的宽度及相邻条形磁场区域的间距均为 $d=0.5\Um$，现有一边长 $l=0.2\Um$、质量 $m=0.1\Ukg$、电阻 $R=0.1\UO$ 的正方形线框 MNOP 以 $v_0=7\Ums$ 的初速从左侧磁场边缘水平进入磁场，求：

\begin{enumerate}
	\item
	线框 MN 边刚进入磁场时受到安培力的大小 $F$；
	
	\item
	线框从开始进入磁场到竖直下落的过程中产生的焦耳热 $Q$；
	
	\item
	线框能穿过的完整条形磁场区域的个数 $n$。
\end{enumerate}

\onepicture[w=10.0cm, align=right]{\includesvg{figs/18}}

%% answer: 
\jdanswer{
\begin{enumerate}
	\item
	$F=2.8\UN$
	
	\item
	$Q=2.45\UJ$
	
	\item
	$n=4$ 个
\end{enumerate}
}
%% memo:
\memoanswer{
（1）线框 MN 边刚开始进入磁场区域时，切割产生的感应电动势为 $E=Blv_0$ ①，感应电流为 $I=\frac{E}{R}$ ②，线框受到的安培力为 $F=BIl$ ③，由①②③解得 $F=2.8\UN$。

（2）设线框竖直下落时，线框下落了 $H$，速度为 $v_H$，由能量守恒得 $mgH+\frac{1}{2}mv_0^2=Q+\frac{1}{2}mv_H^2$，由自由落体规律得 $v_H^2=2gH$，解得 $Q=\frac{1}{2}mv_0^2=2.45\UJ$。

（3）解法一：只有在线框进入和穿出条形磁场区域时，才产生感应电动势，线框部分进入磁场区域 $x$ 时，产生的感应电动势为 $E=Blv$，感应电流为 $I=\frac{E}{R}$，线框受到的安培力为 $F=BIl$，解得 $F=\frac{B^2l^2}{R}v$，在 $t\sim t+\Delta t$ 时间内，由动量定理得 $-F\Delta t=m\Delta v$，求和 $\sum\frac{B^2l^2}{R}v\Delta t=\sum\frac{B^2l^2}{R}\Delta x=mv_0$，解得 $\frac{B^2l^2}{R}x=mv_0$，穿过条形磁场区域的个数为 $n=\frac{x}{2l}$，解得 $n\approx4.4$，可穿过 4 个完整条形磁场区域。

解法二：线框穿过第 1 个条形磁场左边界过程中，产生的平均感应电动势为 $\overline{E}=\frac{Bl^2}{\Delta t}$，平均感应电流为 $\overline{I}=\frac{\overline{E}}{R}$，线框受到的平均安培力为 $\overline{F}=B\overline{I}l$，由动量定理得 $-\overline{F}\Delta t=mv_1-mv_0$，解得 $-\frac{B^2l^3}{R}=mv_1-mv_0$，同理线框穿过第 1 个条形磁场右边界过程中，有 $-\frac{B^2l^3}{R}=mv_1'-mv_1$，所以线框穿过第 1 个条形磁场过程中有 $-\frac{2B^2l^3}{R}=mv_1'-mv_0$，设线框能穿过 $n$ 个条形磁场，则有 $-n\frac{2B^2l^3}{R}=0-mv_0$，解得 $n=\frac{mv_0R}{2B^2l^3}\approx4.4$，可穿过 4 个完整条形磁场区域。
}

\item
%% number: 19
%% paperName: 2007年江苏高考第 19 题 16 分
%% typeId: 6
%% type: 计算题
%% chapter: 功和能
%% point: 动能定理
%% method: 无
%% score: 16
%% degree: 250
%% duplicateId: 0
%% body:
如图所示，一轻绳吊着粗细均匀的棒，棒下端离地面高 $H$，上端套着一个细环。棒和环的质量均为 $m$，相互间最大静摩擦力等于滑动摩擦力 $kmg$（$k>1$）。断开轻绳，棒和环自由下落。假设棒足够长，与地面发生碰撞时，触地时间极短，无动能损失。棒在整个运动过程中始终保持竖直，空气阻力不计。求：

\begin{enumerate}
	\item
	棒第一次与地面碰撞弹起上升过程中，环的加速度；
	
	\item
	从断开轻绳到棒与地面第二次碰撞的瞬间，棒运动的路程 $s$；
	
	\item
	从断开轻绳到棒和环都静止，摩擦力对环及棒做的总功 $W$。
\end{enumerate}

\onepicture[w=2.4cm, align=right]{\includesvg{figs/19}}

%% answer: 
\jdanswer{
\begin{enumerate}
	\item
	$a_{\text{环}}=(k-1)g$，方向竖直向上
	
	\item
	$s=\frac{k+3}{k+1}H$
	
	\item
	$W=-\frac{2kmgH}{k-1}$
\end{enumerate}
}
%% memo:
\memoanswer{
（1）设棒第一次上升过程中，环的加速度为 $a_{\text{环}}$，环受合力为 $F_{\text{环}}=kmg-mg$ ①，由牛顿第二定律得 $F_{\text{环}}=ma_{\text{环}}$ ②，由①②得 $a_{\text{环}}=(k-1)g$，方向竖直向上。

（2）设以地面为零势能面，向上为正方向，棒第一次落地的速度大小为 $v_1$，由机械能守恒得 $\frac{1}{2}(2m)v_1^2=2mgH$，解得 $v_1=\sqrt{2gH}$，设棒弹起后的加速度为 $a_{\text{棒}}$，由牛顿第二定律得 $a_{\text{棒}}=-(k+1)g$，棒第一次弹起的最大高度为 $H_1=-\frac{v_1^2}{2a_{\text{棒}}}$，解得 $H_1=\frac{H}{k+1}$，棒运动的路程为 $s=H+2H_1=\frac{k+3}{k+1}H$。

（3）解法一：棒第一次弹起经过的时间为 $t_1$，与环达到相同速度 $v_1'$，环的速度为 $v_1'=-v_1+a_{\text{环}}t_1$，棒的速度为 $v_1'=v_1+a_{\text{棒}}t_1$，环的位移为 $h_{\text{环}1}=-v_1t_1+\frac{1}{2}a_{\text{环}}t_1^2$，棒的位移为 $h_{\text{棒}1}=v_1t_1+\frac{1}{2}a_{\text{棒}}t_1^2$，又 $x_1=h_{\text{环}1}-h_{\text{棒}1}$，解得 $x_1=-\frac{2H}{k}$，棒环一起下落至地时，有 $v_2^2-v_1^2=2gh_{\text{棒}1}$，解得 $v_2=\sqrt{\frac{2gH}{k}}$，同理，环第二次相对棒的位移为 $x_2=-\frac{2H}{k^2}$，$\ldots$，同理，环第 $n$ 次相对棒的位移为 $x_n=-\frac{2H}{k^n}$，环相对棒的总位移为 $x=x_1+x_2+\ldots+x_n$，摩擦力对棒及环做的总功为 $W=kmgx$，解得 $W=-\frac{2kmgH}{k-1}$。

解法二：设环相对棒滑动距离为 $l$，由能量守恒得 $mgH+mg(H+l)=kmgl$，摩擦力对棒及环做的总功为 $W=-kmgl$，联立解得 $W=-\frac{2kmgH}{k-1}$。
}

\end{enumerate}

\end{document}
